\chapter{Ni dieux, ni maîtres, ni DRM}

\begin{quote}
    \textit{ Ceci n’est ni un parti, ni un programme, ni un slogan, ni une méthode, ni quoi que ce soit d’autre qu’une idée à développer. Postulat : Chaque action sur le monde, de la plus petite à la plus grande, mérite d’être observée méticuleusement et ses conséquences envisagées à l’échelle globale, lointaine, holistique. }\\

    Anonyme, \textit{Manifeste pour une révolution liquide}, 2019
\end{quote}


Dans ce deuxième chapitre, nous étudierons la perception du monde des individus qui composent les communautés patrimoniales pirates. Nous nous intéresserons d'abord aux sentiments de sédition et de rébellion des individus et des communautés face à l'ordre des choses. Ensuite, nous nous concentrerons sur les émotions de peur, d'inquiétude et d'incertitude face aux abysses numériques. Enfin, nous creuserons les parallèles entre la vie culturelle en ligne et les engagements personnels dans le monde physique. Nous nous inscrirons ici dans le travail sur les émotions patrimoniales de Daniel Fabre \footnote{\cite{fabre_emotions_2013}} et les travaux qui enrichissent ce travail, comme ceux d'Emmanuelle Bermès en l'introduisant dans des contextes de communautés patrimoniales numériques \footnote{\cite{bermes_lecran_2024}}. Nous nous écarterons également du pur travail de restitution d'entretiens que nous avons fournis dans le chapitre précédent. Nous nous plongerons dans les traces de plusieurs communautés patrimoniales afin d'affiner ce lien entre émotions individuelles et ces communautés.

\section{Les émotions de la clandestinité}

L'une des plateformes les plus intéressantes à explorer et observer est sans conteste Reddit. Surnommé  \textit{Le cœur de l'Internet}, ce site d'agrégation de forums permet de trouver des communautés dédiées à n'importe quel sujet. Concernant la piraterie culturelle, Reddit et ses forums dédiés sont d'une accessibilité, d'une transparence et d'une documentation rares et précieuses pour le milieu et ses codes. 

\subsection*{La peur, l'inquiétude et la prise de risque}

Dans ce monde, la transparence et une documentation claire sont l'exception qui confirme la règle, la plupart des services d'accès au piratage culturel demandent un minimum de compétences et des savoirs prérequis. Cette opacité totale, doublée de l'inquiétude de télécharger des fichiers caviardés, d'être sur les radars des autorités sanctionnant le piratage, fait que le vécu des pirates ressemble bien plus à une traversée dans l'inconnu qu'à une chasse au trésor. Dans les verbatims de nos enquêtés, l'émotion présente à chaque étape du processus d'exploration, d'acquisition de compétences et de découverte de communautés pirates est la peur ou la crainte. 

\begin{quote}
    A - \textit{Je pense qu'il y a eu une partie de peur au début, quand c'est moi qui m'y suis mis à le faire. Quand mon père le faisait, je regardais derrière son écran, et ça avait l'air trop bien et super facile. Quand c'est moi qui ai dû aller sur des sites un peu shady, avec les extensions, c'est pas que .fr, .com, .org, mais des trucs un peu bizarres, qui sont en fait normal mais sur le coup t'es pas habitué, comment ça .xyz ? C'est quoi cette merde .me ? [...] Je pense que la première étape c'était un peu la peur, un peu l'excitation de ses découvertes [...] de se rendre compte que t'as une sorte de monde, enfin pas un monde caché ce serait un peu abusé, mais de découvrir au fur et à mesure tout ce qui possible de faire avec pas grand chose.}
\end{quote}

\begin{quote}
    T - \textit{ J'avais un fond de crainte et je n'étais pas complètement renseigné non plus, mais j'ai toujours été ce qu'on appelle le mauvais utilisateur dans le torrent c'est celui qui vient se servir et qui n'aide personne [...] parce que c'est comme ça que tu te fais avoir, c'est quand toi tu émets des fichiers, la plupart du temps pour le torrent. [...] Je triche, parce que je fonctionne aussi avec des sites qui demandent à ce que tu équilibres ton ratio, mais tu as des manières de tricher, t'as des manières de fausser l'envoi de certains fichiers et ça compense à chaque fois ce que tu télécharges. [...] Je jour contre les  règles, mais la logique derrière c'est pas d'être le vilain de l'histoire, c'est de ne pas me faire avoir.}
\end{quote}

La source principale d'inquiétude et de peur chez les pirates sont les autorités judiciaires. Chez nos enquêtés, la période d'activité d' Hadopi a laissé un mélange de crainte et de défiance. La majorité d'entre eux a reçu ou a connu une personne ayant reçu une des lettres et des mails de la haute autorité. Si aujourd'hui les lettres Hadopi sont une histoire rigolote pour les pirates francophones, les mouvements étatiques pour le contrôle et la surveillance du web, en particulier pour le piratage, ont produit et contribué à alimenter de réelles inquiétudes dans les communautés. 

\begin{quote}

A - \textit{C'est aussi le moment où j'ai découvert le mot-clé Hadopi, quand j'ai reçu une lettre d'Hadopi, en fait, chez moi. [...] Ouais, j'ai reçu une lettre Hadopi, mais ouais je m'en souviens, c'était un 31 octobre je crois la date même. Alors la question c'était est-ce que c'était moi ou est-ce que c'était le beau-père de l'époque ? On ne sait pas, mais l'adresse IP a été fag, on a reçu une lettre d'Hadopi, et c'est à ce moment-là que mon père m'a payé l'abonnement à la seedbox, comme ça je pouvais télécharger autant que je voulais à l'extérieur, et le principe des torrents étant quand même de télécharger pour rendre ensuite. La seedbox était constamment ouverte, donc sur un serveur, que je pouvais seed constamment, les trucs que je téléchargeais pour pouvoir les envoyer aux autres.}
    
\end{quote}

\subsection*{Trouver et utiliser un service clandestin}

Au-delà de la réflexion tautologique  \og les communautés pirates sont difficiles à découvrir car elles sont illégales et donc se cachent \fg, il est intéressant d'essayer de comprendre le fonctionnement et les organisations mises en place par ces communautés pour survivre et mettre en place des services communautaires. Dans notre introduction, nous avions invoqué le terme \textit{d'organisation de guérilla anarchique} pour décrire l'organisation de mise à disposition des services de piraterie et des lieux de rencontre et d'échange des communautés patrimoniales pirates. Découpons ce terme en deux éléments : \og organisation de guérilla \fg et \fg organisation horizontale \fg.

\begin{itemize}
    \item \textit{L'organisation de guérilla} : Le terme d'organisation de guérilla est une analyse produite par Martin Paul Eve dans son étude sur l'organisation minimaliste des bibliothèques pirates qu'il étudie comme Sci-hub et Libgen \footnote{Martin Paul Eve, \textit{Lessons from the Library: Extreme Minimalist Scaling at Pirate Ebook Platforms}, dans Digital Humanities Quarterly, Providence, vol n°16, n°2, 2022}. 

    \begin{quote}
         \textit{\og Ces archives, qui violent les droits d'auteur, contournent les murs payants et offrent un accès à tous les visiteurs. Bien qu'ils soient souvent du mauvais côté de la loi, les opérateurs de bibliothèques fantômes conçoivent leurs sites en termes éthiques. En définissant leur banditisme en termes de Robin des Bois, ces sites croient qu'ils volent aux riches pour donner aux pauvres de la recherche. Ce qui est important pour les thèmes de ce numéro spécial, c'est que je pense que de telles archives peuvent également être comprises en termes de minimalisme informatique. Comme je le détaillerai plus loin, divers principes de conception technique de ces archives réduisent les barrières sociales à la participation.\fg }

         Martin Paul Eve, traduit, paragraphe 7
    \end{quote}

    Cette analyse peut être mis en relation avec le \textit{Guerilla Open Access Manifesto} attribué au jeune Aaron Swartz, qui a popularisé le terme de guérilla pour décrire à la fois le mode d'organisation des communautés pirates, et la considération de lutte contre une hégémonie mercantiliste de la publication de la connaissance et de la culture.

    \item \textit{L'organisation horizontale} : Le lexique de l'horizontalité fait référence à l'organisation autogestionnaire et l'idée de libre association permettant l'émergence d'une organisation fluide permettant la survit et la réplication des lieux d'échanges des communautés pirates. L'important est la communauté plutôt que le lieu où elle se réunit, n'importe qui peut se revendiquer de la communauté du moment qu'il obéit aux règles parfois tacites, parfois écrites de cette communauté. 
\end{itemize}

L'organisation et le mode de fonctionnement des bibliothèques mis en ligne par les communautés patrimoniales pirates ont un double objectif : être à la fois simples et minimalistes tout en offrant un minimum de barrières à l'entrée et des fonctionnalités avancées et utiles aux utilisateurices. Le minimalisme provient à la fois d'un objectif de réduction des coûts de développement, de maintenance et d'hébergement des bibliothèques pirates, mais également d'offrir le moins de surface visible possible. Les bibliothèques pirates comme Sci-Hub, LibGen, Zlib se présentent souvent comme une simple barre de recherche sur un fond blanc, sans besoin d'authentification ou d'inscription utilisateur. Ce minimalisme est également un mécanisme de défense, le but d'une bibliothèque pirate est d'être le plus camouflé possible, tout en offrant la plus grande qualité et quantité de services possibles pour être partagé le plus globalement possible. Un site relativement simple sera plus facile à remettre en ligne si un domaine, du matériel, des individus sont éliminés ou emprisonnés par les autorités judiciaires. Nous reviendrons sur l'organisation des dépôts des communautés patrimoniales pirates dans le chapitre suivant, mais l'important est la différence de philosophie qui habite le design et l'utilisation des bibliothèques pirates comparée aux applications et aux services légaux. Trouver et utiliser ces outils demande des codes sociaux, des habitudes différentes, et les premiers pas dans cet univers sont souvent anxiogènes et peuvent faire tourner les talons aux individus les moins bien accompagnés.

\begin{quote}
    E - \textit{Mais plus c'est difficile d'accéder, plus je vais avoir peur d'aller fouiner, finalement. D'aller chiner des trucs, parce que je me dis que c'est caché pour une raison, je ne sais pas. Mais sinon, au-delà de la peur, c'est vraiment du soulagement, de l'euphorie ? Euphorie c'est peut-être beaucoup, j'avais besoin d'un film, je l'ai trouvé. Ça fait un peu un saut de dopamine. [...] Maintenant je n'ai plus trop peur, malgré ... Le souci des sites de streaming, c'est qu'il y a tout le temps des pubs. Des fois, j'ai un peu peur d'avoir un virus qui s'installe. Ça, c'est toujours un petit peu ... ça je n'aime pas trop.}
\end{quote}

\section{Du partage des émotions aux communautés pirates }

Dans son livre \textit{Du numérique aux émotions}, Emmanuelle Bermès adapte la typologie des émotions patrimoniales de Daniel Fabre au contexte des communautés se réunissant autour d'objets numériques. L'émotion mise en avant pour comprendre la mentalité des mouvements \og hacktivistes \fg qui mettent en avant le partage radical de l'information scientifique est la sédition. L'envie de rompre avec un système considéré comme injuste et dysfonctionnel permet de réunir des individus autour de l'idée de créer l'alternative au sein de l'infrastructure, par la mise en place d'outils, l'autogestion et une pratique révolutionnaire (ne pas jouer selon les règles du jeu) ; plutôt que d'attendre la lente évolution de la superstructure.

\subsection*{Du pirate à l'archiviste rebelle : de la consommation à une pratique politique}

Dans les entrevues avec nos enquêtés, la question du rapport politique de la pratique du piratage n'a pas été une piste ayant produit des réponses très concluantes. Nous avions déjà cité l'exemple d'Hilaire, de ses considérations éthiques et de sa découverte de la culture du libre dans le chapitre précédent. Dans le cas de Théo, la piraterie culturelle s'inscrit dans une vision globale d'un idéal d'accès à la culture ; cependant, il existe un fossé entre une vision du monde et un engagement politique personnel que Théo se refuse à franchir. De son côté, Zakary a produit des réflexions sur l'éthique de sa pratique, considérant l'organisation économique du monde des auteurs de mangas. 

\begin{quote}
    Z - \textit{Si je dois payer un élément, je dois être sûr en tant qu'acheteur, en tant que consommateur, que de mon côté je vais avoir un produit de qualité... Et un produit, du coup, assez éthique. Je dois juste être sur de la qualité, et de l'autre côté de l'artiste, il y a un traitement éthique. Il faut que l'artiste soit payé convenablement, que l'argent que je lui fournis soit bien touché, etc. Et tout simplement j'ai l'impression que moi, en tant que pirate, ce contrat a été rompu par toutes les plateformes qui proposent des alternatives légales : Netflix, Amazon, Crunchyroll, ADN, Manga Plus. Mais c'est toujours un peu critiqué dans la communauté.}
\end{quote}

Cependant, aucun de nos enquêtés ne revendique la piraterie culturelle comme la base d'une identité politique, même si certains d'entre eux se considèrent comme des personnes politisées et militantes. Dans différentes communautés patrimoniales, la sédition et la consommation pirate revêtent une couleur politique appelant à remettre en cause l'organisation capitaliste de certaines ressources. L'exemple réapparaissant régulièrement est la disparition d'un contenu culturel de la vente physique et des services de location. Certains mouvements se sont structurés autour de la lutte à offrir une vie aux objets culturels après l'exploitation commerciale, pour que la culture ne soit pas à la main des ayants droit. Par exemple, le terme d'\textit{abandonware} correspond à un logiciel qui n'est plus maintenu ni commercialisé par ses ayants droits. D'une certaine manière, l'horizon de tout objet numérique est de tomber dans l'oubli et de disparaître ou de devenir l'équivalent d'un abandonware, un objet culturel orphelin. La conscience de l'existence de la perte d'accès aux contenus culturels existe, sans réellement porter de temps. Au-delà de la question d'accès à la science et à l'éducation, à la barrière monétaire d'accès à la culture, la peur de l'oubli collectif des contenus culturels est l'une des motivations les plus profondes des communautés patrimoniales pirates.

\begin{quote}
\textit{    > Check Amazon, Netflix, Amazon, Disney, Max, Apple, etc.}

\textit{    > No where to find Yukio Mashima films
}

\textit{    > No where to find Spaceballs
}

\textit{    > No where to find Amadeus
}

\textit{    > No where to find Samurai (1954)
}

\textit{    > No where to find Jhin-Rho 
}

\textit{    They are just not "available", Why ?
}

    Anonymous n°210962625, \og My country has a Kino embargo \fg, sur 4chan, le 28/05/2025, via le post de u/Vor\_Mor, \og  No kino movie \fg, sur r/greentext, 28/05/2025.
\end{quote}

La base philosophique de la piraterie culturelle des individus et des communautés qu'ils composent est clairement d'inspiration libertaire. Cependant, il serait erroné de considérer que les forums d'échange de films sont peuplés en majorité de militants anarchistes. %A compléter


\subsection*{Les espaces numériques de discussions}

Les forums dédiés à la piraterie culturelle agissent de la même manière qu'une agora, où certains événements font office de catalyseurs amplifiant les émotions individuelles. Ces espaces de discussion et de partage deviennent le berceau d'actions concrètes. Les subreddits r/Piracy et spécifiquement r/DataHoarder sont des sources fascinantes sur les communautés patrimoniales numériques. Premièrement, ce sont des forums ouverts ; ces dernières années, une grande partie des communautés de fans, de travaux bénévoles sur des projets open source, des wikis, de communautés patrimoniales pirates a migré sur des logiciels de discussion communautaire semi-fermés comme Discord. C'est un outil fabuleux pour offrir un espace d'échange et construire une communauté ; cependant, sa nature d'espace d'échange privé a pour conséquence de supprimer du web indexable sur un moteur de recherche une grande partie des ressources et des échanges qui documentent l'existence et l'évolution de ces communautés. Deuxièmement, ces subreddits possèdent des outils de recherche intégrés comme le tri par ordre chronologique, par popularité des posts et des réponses, des controverses (les posts ayant le plus de réponses mais peu de \og upvotes \fg, l'équivalent des likes sur reddit). Ceci nous permet de suivre les discussions, les débats et les intérêts des communautés patrimoniales pirates.

\begin{quote}

     \textit{I think there will be a Jan 6 situation where this will get wiped off the internet, are there any current efforts to archive footage and images from this current ongoing event? If not I'd think that's something that should be payed attention to at the moment.} 

     u/awolfwearingabanana, \og Any attempts to archive the current LA protests? \fg, sur r/DataHoarder, le 09/06/2025
\end{quote}

Ce poste de u/awolfwearingabanana, sur l'archivage amateur de la documentation et des vidéos des émeutes ayant cours à Los Angeles en juin 2025, nous permet de plonger directement dans les discussions et les débats qui animent la communauté des \textit{archivistes rebelles} de r/DataHoarder.  Très rapidement, les premiers commentaires des utilisateurs les plus actifs font un parallèle avec les émeutes du Capitole du 6 janvier 2021, et partagent les liens torrent vers les dépôts des utilisateurs qui se sont donnés comme mission d'archiver la documentation de ces émeutes. L'expérience de la communauté des archivistes d'images et de vidéos de manifestation et les différentes méthodologies de sélection et d'extraction de vidéos sont discutées sous ce poste. D'autres creusent la distinction entre l'amassement de données, l'archivage et la mise à disposition du public dont ces données auraient besoin. L'initiative ne fait cependant pas l'unanimité, et une partie des membres du forum critique la propension à une forme de paranoïa archivistique, un complotisme de la mémoire. 

\begin{quote}

\textit{My suggestion would be to download content you find. Assume others are not saving anything, then share your collection.
There's a 1TB torrent of Jan6 content which I believe was a compilation of files different people downloaded and then shared here.
(edit: I'm not from the US or saving any content. Just suggesting what you can do and pointing out what others done before.)}

u/P03tt, en réponse à u/awolfwearingabanana sur r/DataHoarder.\\

\textit{Archival is the first step, and what a lot of people do, but I would argue it's just the first step. If you're looking at preserving data for the public good, you should look at making those archives available, then curate and index them. If you do create an archival site, feel free to let us know and we can add it to the https://datahoarding.org/ index. We already have a lot of Jan6 archives on there. }

u/shimoheihei2, en réponse à u/awolfwearingabanana sur r/DataHoarder.\\

\textit{Jan 6 situation where this will get whipped off the internet
Except for the fact that this never happened? Everyone knows what Jan 6 is, all the videos and photos are online and easily accessible.
I swear redditors live in this weird meta-world where they believe reddit is simultaneously some kind of safe harbor for talking about events the rest of the world would get "removed" for. If "they" (whoever you believe wants to "whip" things of the internet) is able to actually do so, why do you believe this post is up? If what you believe was real you'd be taken into a black site and "questioned" while your post got removed the second it hit a server. payed .Where's the bot? Oh wait, no, I get it. You're on a boat stuck with very slow satellite internet and believe everything is being censored because it won't load on your shitty connection?}

u/IronCraftMan, en réponse à u/awolfwearingabanana sur r/DataHoarder.\\

\end{quote}

\section{Conclusion}

L'horizon mental et émotionnel des individus qui piratent peut être étudié sous le prisme des théories de la sociologie de la déviance. La piraterie culturelle, avec les définitions de la sociologie, peut être considérée comme un fait normal \footnote{Albert Ogien, \textit{Le normal et le pathologique}, dans \og Sociologie de la déviance \fg, p. 29-38, Presses Universitaires de France, 2012}, que l'on retrouve uniformément dans des contextes séparés. L'élément de réflexion qui rend l'étude de la piraterie culturelle très intéressant est la jonction entre un acte de déviance (je refuse de payer pour un objet numérique payant) et ses conséquences (peur, clandestinité, prise de risque), une revendication politique et culturelle (conservation, éducation, rapport aux institutions et entreprises ayant les droits d'exploitation), et un aspect communautaire spécifique à la dimension patrimoniale (forums, organisation de partage et de collecte, bibliothèques pirates). La massification de la capacité de copie des individus entre en contradiction avec le système de distribution culturelle et les règles qui l'entourent. Le sentiment de dépossession face au patrimoine numérique pousse vers la sédition et le piratage comme catalyseur d'alternatives. Cet espace clandestin demande un apprentissage social et culturel pour pouvoir être navigué, mais il renforce les liens communautaires et permet aux individus de se sentir maîtres de leur consommation culturelle en forgeant de nouvelles compétences au cours de leurs carrières pirates. L'apprentissage d'une culture de la piraterie culturelle permet aux individus de se regrouper en communautés, avec comme moteur les émotions générées par le rapport de la société aux objets culturels immatériels. Ces considérations pour des objets culturels rentrent dans le domaine du politique, même si les individus ne sont pas forcément des militants politiques. 

